\section{Variations, order, and the Ising entropy}
\label{var:section}

Let $f\in C_c^\infty(\mathbb R)$. The upper bound will use two
properties of the variational formula: increasing the tail integrals of
an overlap path decreases $V_f$, and finite step trials with top value
below one suffice for the infimum. The lower bound needs a different
property: concavity of $\ell$ turns its first variation into a global
supporting inequality, which will identify the entropy supremum.
We establish these properties from the recursions of
Definition~\ref{def:recursions}.

Recall their different endpoints:
the pattern recursion has $q_{-1}=0$, $q_k=1$, and coefficients
$\zeta_0=0,\ldots,\zeta_k=1$, whereas the field recursion has
terminal value $\log\cosh(x)-h_{k-1}/2$ and stops at level $k-1$.
Inserting a partition point without changing a step value inserts an
identity transform of variance zero. Consequently either recursion may
be computed on any common refinement of the partitions in use.

\subsection{First variations and the pattern tail order}

For either recursion write $s_i$ for its variance increment and define
the probability kernel
\begin{equation}
\label{var:kernel}
 K_i(x,dy)=
 \exp\{\zeta_i(U_i(y)-U_{i-1}(x))\}
 \mathcal N(x,s_i)(dy).
\end{equation}
Here $s_i=q_i-q_{i-1}$, $0\le i\le k$, in the pattern
recursion, and $s_i=h_i-h_{i-1}$, $0\le i\le k-1$, in
the field recursion. The normalization follows directly from the
definition of $\mathcal T_{s_i,\zeta_i}$; when $\zeta_i=0$
the exponential is one. A zero-variance kernel is the point mass at
its starting point. Let $Z_{-1}=0$ and, successively, sample $Z_i$
from $K_i(Z_{i-1},\cdot)$.

Two elementary derivative identities will be used repeatedly. For a
smooth terminal function $U$, put $W=\mathcal T_{s,d}U$ and let
$K$ denote the corresponding tilted Gaussian kernel. Differentiation
in $x$ and in a perturbation $U+\varepsilon A$ gives
\begin{equation}
\label{var:transform-derivatives}
 W'=K U',\qquad
 \left.\frac{d}{d\varepsilon}
 \mathcal T_{s,d}(U+\varepsilon A)\right|_{\varepsilon=0}=K A.
\end{equation}
The Gaussian heat equation gives
\begin{equation}
\label{var:generator}
 \partial_s W
 =\frac12\{W''+d(W')^2\}
 =\frac12 K\{U''+d(U')^2\}.
\end{equation}
For $d>0$, these identities follow by differentiating
$e^{dW}=\mathbb E e^{dU(x+\sqrt sG)}$; for $d=0$ they
are the ordinary heat-semigroup identities. In particular,
$U_i'(Z_i)$ is a martingale along the tilted chain, since
$K_iU_i'=U_{i-1}'$.

\begin{lemma}[First variations]
\label{var:variation}
On a fixed partition with $w_i=\zeta_{i+1}-\zeta_i$, the first
variations along every feasible line segment of step paths are
\begin{align}
\label{var:pattern-variation}
 \frac{d}{da}V_f(q+a\dot q)
 &=-\frac12\sum_{i=0}^{k-1}w_i\dot q_i\,
   \mathbb E[U_i'(Z_i)^2],\\
\label{var:field-variation}
 \frac{d}{da}\ell(h+a\dot h)
 &=-\frac12\sum_{i=0}^{k-1}w_i\dot h_i\,
   \mathbb E[U_i'(Z_i)^2].
\end{align}
The functions and chains on the right are computed at the current
path. Endpoint derivatives are understood from within the segment.
The squared-derivative expectations on the right are continuous on the
closed finite-dimensional parameter cones. In both recursions these
expectations are nonnegative and nondecreasing in $i$. In the field
recursion, define
\begin{equation}
\label{var:rho}
 \rho_h(u)=\rho_i:=\mathbb E[U_i'(Z_i)^2],
 \qquad \zeta_i<u<\zeta_{i+1}.
\end{equation}
Then $0\le\rho_h\le1$, and subdivision does not change this path.
\end{lemma}

\begin{proof}
First suppose all variance increments being differentiated are
strictly positive. Moving $q_i$ lengthens the transition ending at
level $i$ and shortens the transition leaving that level. By
\eqref{var:generator}, the incoming derivative, expressed at level
$i$ before propagation through the preceding kernels, is
$\{U_i''+\zeta_i(U_i')^2\}/2$. The outgoing derivative is
$-\{U_i''+\zeta_{i+1}(U_i')^2\}/2$. Their sum is
$-w_i(U_i')^2/2$. Formula~\eqref{var:transform-derivatives}
propagates this sum through the kernels ending at levels
$i,i-1,\ldots,0$, giving
$\partial_{q_i}V_f=-w_i\mathbb E[U_i'(Z_i)^2]/2$.

The same cancellation applies to $h_i$ for $i<k-1$.
At the last field level the incoming derivative and the derivative
of the subtraction $-h_{k-1}/2$ combine to give
\[
 \frac12\{(\log\cosh)''+
            \zeta_{k-1}((\log\cosh)')^2-1\}
 =-\frac{1-\zeta_{k-1}}2\tanh^2.
\]
This is the same formula with $w_{k-1}=1-\zeta_{k-1}$.
Linearity yields both displayed variations.

For completeness, the differentiations remain valid at degenerate
increments. A Gaussian transform of either terminal function is
smooth in the spatial variable and has bounded derivatives of every
positive order. This follows inductively by differentiating the
tilted expectation: the resulting terms are finite sums of tilted
moments of products of bounded derivatives of the next-level
function. The functions themselves are bounded in the pattern
case and have at most linear growth in the field case. On every
compact variance-parameter set, the Gaussian exponential moments
therefore dominate all terms in these formulas. Coupling the
increments as $\sqrt{s_i}G_i$ shows continuity as any $s_i$
decreases to zero, both for recursion values and for the tilted
moments that appear above. The heat-equation identity also has
this continuous extension because its spatial derivatives do.
Approximate the endpoints of a feasible segment by interior points
on a common partition and integrate the already proved derivative
formula along the approximating segments. Dominated convergence
gives the same integral formula on the original segment and its
one-sided endpoint derivatives.

Finally, the derivative martingale and conditional Jensen inequality
show that its second moments increase with the level. The terminal
derivative bounds are $\|f'\|_\infty$ and $1$, respectively;
the same bounds hold at every earlier level by
\eqref{var:transform-derivatives}. An inserted zero-variance
transition repeats the same derivative and chain value, so it
does not change $\rho_h$.
\end{proof}

\begin{proposition}[Continuity and tail order]
\label{var:continuity-order}
For step paths on arbitrary partitions,
\begin{align}
\label{var:lipschitz}
 |V_f(p)-V_f(q)|&\le\tfrac12\|f'\|_\infty^2\|p-q\|_{L^1},\\
 |\ell(h')-\ell(h)|&\le\tfrac12\|h'-h\|_{L^1}.
\end{align}
Thus $V_f$ extends uniquely and continuously to all nondecreasing
$[0,1]$-valued paths, and $\ell$ extends to all bounded
nonnegative nondecreasing paths. Moreover,
\begin{equation}
\label{var:tail-order}
 \int_s^1p(u)\,du\ge\int_s^1q(u)\,du
 \quad\hbox{for every }s\in[0,1]
 \quad\Longrightarrow\quad V_f(p)\le V_f(q).
\end{equation}
\end{proposition}

\begin{proof}
Refine the two partitions and integrate
Lemma~\ref{var:variation} along the segment joining the paths.
The terminal derivative bounds give \eqref{var:lipschitz}.
Step paths are dense in the indicated classes in $L^1$, which
gives the extensions.

For the order assertion on a common partition, set
$a_i=w_i(p_i-q_i)$ and $A_j=\sum_{i=j}^{k-1}a_i$.
The hypothesis gives $A_j\ge0$. For every nonnegative
nondecreasing sequence $c_i$,
\begin{equation}
\label{var:summation-parts}
 \sum_{i=0}^{k-1}a_ic_i
 =c_0 A_0+\sum_{j=1}^{k-1}(c_j-c_{j-1})A_j\ge0.
\end{equation}
At each point of the segment from $q$ to $p$, take
$c_i=\mathbb E[U_i'(Z_i)^2]$ and apply
\eqref{var:pattern-variation}. The derivative along the segment
is nonpositive, proving the claim for steps.

For general paths, replace each by its average on each interval of
a common sequence of successively finer uniform partitions.
These averages are nondecreasing and converge in $L^1$ to the
original paths. Their tail-integral differences at partition
boundaries equal the original differences. Between boundaries
the differences are affine and hence nonnegative as well.
Apply the step result and then \eqref{var:lipschitz}.
\end{proof}

\begin{corollary}[Finite variational value and reduction to finite trials]
\label{var:step-reduction}
The entropy in Definition~\ref{def:entropy} is nonnegative,
nondecreasing in pointwise path order, and satisfies
$S_{\mathrm I}(0)=0$. The variational infimum is finite, and
\begin{equation}
\label{var:step-infimum}
 \inf_q\{\alpha V_f(q)+S_{\mathrm I}(q)\}
 =\inf_{\substack{q\text{ step}\\q_{k-1}<1}}
       \{\alpha V_f(q)+S_{\mathrm I}(q)\}.
\end{equation}
\end{corollary}

\begin{proof}
The trial $h=0$ gives $S_{\mathrm I}(q)\ge0$. Each expression
in its defining supremum is nondecreasing in $q$ because $h\ge0$.
Also $\ell(0)=0$ and \eqref{var:field-variation}, integrated
along $a\mapsto ah$, gives $\ell(h)\le0$, whence
$S_{\mathrm I}(0)=0$. Every transform preserves the range of a
bounded terminal function, so $|V_f(q)|\le\|f\|_\infty$.
In fact, the zero path has
$V_f(0)=\log\mathbb E e^{f(G)}$. These observations bound the
infimum below by $-\alpha\|f\|_\infty$ and provide a finite
upper bound using $q=0$.

For any quantile $q$, let $\underline q_m$ on each interval of
the uniform $m$-partition be $(1-1/m)$ times the essential
infimum of $q$ on that interval. Then
$0\le\underline q_m\le q$ almost everywhere,
$\sup\underline q_m<1$, and
$\|\underline q_m-q\|_{L^1}\to0$. The last assertion follows,
for example, from convergence at every continuity point of a
monotone representative, followed by bounded convergence.
Consequently $S_{\mathrm I}(\underline q_m)\le S_{\mathrm I}(q)$,
while $V_f(\underline q_m)\to V_f(q)$.
Taking a limit superior for each finite-entropy trial $q$ proves
the nontrivial inequality in \eqref{var:step-infimum}; the other
inequality follows from inclusion of the trial classes. No
continuity assertion for the extended-valued entropy is needed.
\end{proof}

\subsection{Concavity of the field recursion}

The supporting inequality below is the property of the centered Ising
spin law used in the lower bound. Its concavity statement is the
path-convexity theorem of Chen, Issa, and
Mourrat~\cite[Propositions~2.2 and~2.4]{chenissamourrat2026}: with their
notation and covariance normalization, $\ell(h)=-\psi_{\circ}(h/2)$.
Ho~\cite[Theorem~1.8]{ho2026concavity} proves the same path-convexity
statement through a concavity theorem for the time-changed Parisi PDE.
This is convexity along linear interpolations of quantile paths, not
along affine interpolations of their probability laws. The proof below
uses the Hessian criterion of Chen, Issa, and Mourrat: the negative Hessian has nonpositive
off-diagonal entries and nonnegative row sums. We derive the
off-diagonal signs directly from folded Gaussian kernels and include
the extension to repeated or zero field levels. The preservation of
even functions monotone in absolute value by these transitions is
classical; see Panchenko~\cite[Lemma~2(d)]{panchenko2005question}.
The next lemma gives a direct proof in the form needed here.

\begin{lemma}[Order preservation by a folded transition]
\label{var:folded-kernel}
Suppose $U:\mathbb R\to\mathbb R$ is even and convex. Let $K$ be
its tilted Gaussian kernel with variance $s\ge0$ and coefficient
$d\ge0$, and assume its normalizer
$\mathbb E e^{dU(x+\sqrt sG)}$ is finite for every $x\in\mathbb R$.
If $A$ is bounded, even, and nondecreasing in $|x|$, then $KA$
has these properties as well. The same statement with
``nonincreasing'' holds with both monotonicities reversed.
\end{lemma}

\begin{proof}
Evenness follows by reflection, and $s=0$ is immediate. For
$s>0$, the law of the next absolute value $y\ge0$ when the
starting point is $x\ge0$ has density proportional to
\[
 e^{-y^2/(2s)}\cosh(xy/s)e^{dU(y)}.
\]
For $x_2\ge x_1\ge0$, the ratio of these normalized densities
is a positive constant times
$\cosh(x_2y/s)/\cosh(x_1y/s)$. Its logarithmic derivative in
$y$ is
\[
 \frac{x_2}{s}\tanh(x_2y/s)
 -\frac{x_1}{s}\tanh(x_1y/s)\ge0,
\]
because $x\mapsto x\tanh(xy/s)$ is nondecreasing on
$[0,\infty)$. Thus the density ratio $L(y)$ is increasing.
If $Y,Y'$ are independent with the law corresponding to $x_1$,
then
\[
 \operatorname{Cov}(A(Y),L(Y))
 =\tfrac12\mathbb E[(A(Y)-A(Y'))(L(Y)-L(Y'))]\ge0.
\]
Since $\mathbb E L(Y)=1$, this proves
$K A(x_2)\ge K A(x_1)$. Negating $A$ proves the reversed
statement. The normalizations are included in $L$ and do not
alter its monotonicity.
\end{proof}

\begin{theorem}[Concavity and a supporting field]
\label{var:concavity}
The functional $\ell$ is concave on nonnegative nondecreasing
step paths. For any two such paths, represented on a common
partition,
\begin{equation}
\label{var:support}
 \ell(h')\le\ell(h)
      -\frac12\int_0^1(h'(u)-h(u))\rho_h(u)\,du.
\end{equation}
Concavity also holds for its continuous extension to bounded
nonnegative nondecreasing paths.
\end{theorem}

\begin{proof}
Fix a partition. We first work in its open cone
$0<h_0<h_1<\cdots<h_{k-1}$, where all parameter derivatives
below exist. Write $m_i=U_i'$ and
$\rho_i=\mathbb E m_i(Z_i)^2$ as in \eqref{var:rho}.
By the first variation, the matrix
\[
 A_{ij}=w_i\partial_{h_j}\rho_i
       =-2\partial_{h_j}\partial_{h_i}\ell
\]
is symmetric. We will prove that its off-diagonal entries are
nonpositive and its row sums are nonnegative. These two properties
make $A$ positive semidefinite, as the decomposition
\eqref{var:psd} below shows. They correspond to increasing one later
field level and to increasing all field levels together, respectively.

Every $U_i$ is even and convex. Evenness is preserved by the
Gaussian transforms. To check convexity when $d>0$, use convexity
of $U$ and H\"older's inequality to obtain, for $0\le a\le1$,
\[
 \mathbb E e^{dU(ax+(1-a)y+\sqrt sG)}
 \le
 (\mathbb E e^{dU(x+\sqrt sG)})^a
 (\mathbb E e^{dU(y+\sqrt sG)})^{1-a}.
\]
Taking logarithms proves convexity of the transform; expectation
preserves it when $d=0$. The terminal subtraction is constant in
$x$ and does not affect the argument. In particular, $m_i$ is
odd, $m_i(x)\ge0$ for $x\ge0$, and $m_i(x)^2$ is even and
nondecreasing in $|x|$. All $U_i$ have at most linear growth, so
the normalizers in Lemma~\ref{var:folded-kernel} are finite.

We first establish the cross-derivative sign
\begin{equation}
\label{var:cross-sign}
 \partial_{h_j}\rho_i\le0\qquad(j>i).
\end{equation}
Repeat the switch calculation from Lemma~\ref{var:variation},
but stop its backwards propagation at the function $U_i(x)$.
This gives
\begin{equation}
\label{var:later-switch}
 \partial_{h_j}U_i(x)
 =-\frac{w_j}{2}
   (K_{i+1}\cdots K_j m_j^2)(x),\qquad j>i.
\end{equation}
Lemma~\ref{var:folded-kernel} says that this function is even
and nonincreasing in $|x|$. Its derivative is therefore
nonpositive for $x>0$ and nonnegative for $x<0$. Since $m_i$
has the opposite signs, the direct derivative of the terminal
test in the expression $\rho_i=K_0\cdots K_i(m_i^2)(0)$ is
pointwise nonpositive:
\[
 \partial_{h_j}(m_i^2)
   =2m_i(\partial_{h_j}U_i)'\le0.
\]

It remains to account for all derivatives of the probability
kernels in this expression. If $l\le i<j$, the variance of
$K_l$ is independent of $h_j$. Put
$A_l=\partial_{h_j}U_l$. Differentiating the normalized kernel
at a fixed starting point gives the score
\begin{equation}
\label{var:kernel-score}
 \partial_{h_j}\log K_l(x,dy)
 =\zeta_l\{A_l(y)-(K_lA_l)(x)\}.
\end{equation}
Indeed the derivative of $U_{l-1}$ in its normalizing factor
equals $K_l A_l$, by
\eqref{var:transform-derivatives}. In the product rule for
$K_0\cdots K_i(m_i^2)$, the test at level $l$ is
$H_l=K_{l+1}\cdots K_i(m_i^2)$, with $H_i=m_i^2$.
It is even and nondecreasing in absolute value by
Lemma~\ref{var:folded-kernel}. On the other hand, $A_l$ is even
and nonincreasing in absolute value by
\eqref{var:later-switch}. Thus differentiating this kernel
contributes, at its starting point,
\[
 \zeta_l\operatorname{Cov}_{K_l(x,\cdot)}(H_l,A_l)\le0.
\]
The sign follows from the same independent-copy covariance
identity as in Lemma~\ref{var:folded-kernel}, now for two
functions of $|y|$ with opposite monotonicities. Propagation
through preceding probability kernels preserves the sign.
This includes $l=0$, whose score is zero because $\zeta_0=0$.
Together with the direct derivative, these are all terms in the
product rule, proving \eqref{var:cross-sign}.

Next increase every field value by the same number $a\ge0$.
All increments after level zero remain unchanged. Each $U_i$,
$i\ge0$, changes only by the additive constant $-a/2$, so its
spatial derivative and all kernels $K_l$ for $l\ge1$ remain
unchanged. The sole change in the chain is that its first,
untilted Gaussian variance increases from $h_0$ to $h_0+a$.
The function
$K_1\cdots K_i(m_i^2)$ is even and nondecreasing in absolute
value. Coupling the initial values as
$\sqrt{h_0+a}G$ therefore proves
\begin{equation}
\label{var:row-sum}
 \sum_{j=0}^{k-1}\partial_{h_j}\rho_i\ge0.
\end{equation}

We can now apply the matrix criterion. By \eqref{var:cross-sign}
and symmetry every off-diagonal entry of $A$ is nonpositive,
and \eqref{var:row-sum} gives nonnegative row sums. For every
$v\in\mathbb R^k$,
\begin{equation}
\label{var:psd}
 v^{\mathsf T}Av
 =\sum_i\left(\sum_jA_{ij}\right)v_i^2
   +\sum_{i<j}(-A_{ij})(v_i-v_j)^2\ge0.
\end{equation}
Thus the Hessian of $\ell$ is negative semidefinite on the
open cone. Restricting $\ell$ to a line segment in that cone
proves concavity and, using its derivative at the starting
point, the inequality \eqref{var:support} there.

For arbitrary field steps $h,h'$ on this partition, replace
$h_i,h_i'$ by $h_i+\varepsilon(i+1)$ and
$h_i'+\varepsilon(i+1)$. Both perturbed paths lie in the open
cone and their difference remains $h'-h$. Let
$\varepsilon\downarrow0$. The value and coefficient continuity
in Lemma~\ref{var:variation} proves both concavity and
\eqref{var:support} on the closed cone. Common refinement
handles distinct partitions. Finally, approximate any two
bounded field paths in $L^1$ by monotone steps on common
partitions and use Proposition~\ref{var:continuity-order} to
pass the concavity inequality to their limits.
\end{proof}

For a step field $h$, the supporting inequality gives the concrete
entropy identity
\[
 S_{\mathrm I}(\rho_h)
 =\ell(h)+\frac12\int_0^1h(u)\rho_h(u)\,du.
\]
Indeed, adding $\frac12\int h'\rho_h$ to \eqref{var:support}
bounds every trial value by the value at $h$, and $h$ itself is
admissible. The nonlinear model must still identify its overlap path
with this field derivative. Section~\ref{cav:section} obtains that
identity from coordinate exchangeability and passes the same supporting
inequality to the bounded field produced by the cavity limit.
